% BCT Superfluid Lattice Model — Volume II: Electroweak I
% Letters L13–L24 | March 2026
% Michel Robert Cabrié | ORCID: 0009-0007-9561-9859
% Truncated Physics Press · Barrys Reef, Victoria, Australia

\documentclass[aps,prl,twocolumn,showpacs,preprintnumbers]{revtex4-2}

\usepackage{amsmath,amssymb,amsfonts}
\usepackage{xcolor}
\usepackage{booktabs}
\usepackage{microtype}
\usepackage{hyperref}

\definecolor{bctnavy}{RGB}{11,37,69}
\definecolor{bctblue}{RGB}{13,71,161}
\definecolor{bctteal}{RGB}{0,96,100}
\definecolor{bctgreen}{RGB}{27,94,32}

\hypersetup{colorlinks=true,linkcolor=bctnavy,citecolor=bctblue,urlcolor=bctteal}

\newcommand{\roct}{r_{\mathrm{oct}}}
\newcommand{\rtet}{r_{\mathrm{tet}}}
\newcommand{\azero}{\alpha_0}
\newcommand{\LQCD}{\Lambda_{\mathrm{QCD}}}

\begin{document}

\preprint{BCT Letters Collection --- Volume II --- March 2026}

\title{BCT Letters Collection --- Volume II: Electroweak~I\\
\large $m_W$, $m_Z$, $m_H$, and the Electroweak Scale}

\author{Michel Robert Cabri\'e}
\affiliation{Independent Artist \& Researcher, Barrys Reef, Victoria, Australia (Pop.\ 28)}
\email{ZeroFreeParameters@gmail.com}
\homepage{https://patreon.com/cw/TheBCTSuperfluidLatticeModel}

\date{March 2026}

\begin{abstract}
Volume~II of the BCT Superfluid Lattice Model programme collects
Letters~13--24, deriving the complete electroweak sector from BCT
vacuum geometry. The W and Z boson masses, the Higgs boson mass,
the electroweak VEV, the Fermi constant, and the weak mixing angle
are all derived from the three geometric inputs with zero free parameters.
Key results: $m_W = 80{,}370$~MeV ($+0.004\%$); $m_Z = 91{,}188$~MeV
($<0.01\%$); $m_H = 125{,}263$~MeV ($+0.010\%$); $v_{\mathrm{BCT}}
= 246.2198$~GeV ($-0.000095\%$). All results use zero free parameters.
Open access CC~BY~4.0.
\end{abstract}

\maketitle

%----------------------------------------------------------------------
\section*{Programme Context}
%----------------------------------------------------------------------

Volume~II derives the electroweak sector of the Standard Model from
BCT geometry. The electroweak scale $v = 246.22$~GeV emerges from
the BCT Gross--Pitaevskii void integral --- not assumed, but derived
from the geometry of the octahedral and tetrahedral voids in the D4
condensate.

The derivation uses the NNLO oct-void metric factor:
\begin{equation}
f_{\mathrm{oct}} = \frac{3}{4}
+ \frac{\alpha_0}{2}(1+\alpha_0)
- \frac{\alpha_2^2}{32}
\end{equation}
where $\alpha_2 = \roct^2/\pi$. This gives $v_{\mathrm{BCT}} = 246.2198$~GeV,
sixty times inside the experimental uncertainty of $246.220 \pm 0.012$~GeV.
All five audited holes in the BCT electroweak derivation are formally
closed in Volume~XX (The Closure Volume).

%----------------------------------------------------------------------
\section{Letter 13 --- The Electroweak Scale from BCT Geometry}

\medskip
The electroweak VEV $v = 246.22$~GeV is derived from the BCT
Gross--Pitaevskii path integral over the octahedral void network.
The derivation identifies $v^2 = 2m_W^2/g_W^2$ where $g_W$ is
the weak coupling, itself determined by the D4 instanton action.
The leading result $v_{\mathrm{BCT}} = 246.2198$~GeV differs from
the PDG value by $-9.5 \times 10^{-5}\%$ --- the most precise
BCT prediction for a dimensionful quantity.

%----------------------------------------------------------------------
\section{Letter 14 --- The W Boson Mass}

\medskip
The W boson mass is derived as:
\begin{equation}
m_W = \frac{v}{2}\,g_W = \frac{v}{2}\sqrt{4\pi\alpha/\sin^2\theta_W}
\end{equation}
where $\sin^2\theta_W = 0.23143$ (derived in Volume~I) and $\alpha$
is the BCT fine structure constant. Result: $m_W = 80{,}370$~MeV,
agreeing with the PDG value of $80{,}367$~MeV to $+0.004\%$
(Prediction~\#108). The BCT prediction precedes the CDF anomaly
and is consistent with the pre-CDF world average.

%----------------------------------------------------------------------
\section{Letter 15 --- The Z Boson Mass}

\medskip
The Z boson mass is derived from the electroweak mixing:
\begin{equation}
m_Z = \frac{m_W}{\cos\theta_W}
= \frac{v}{2}\sqrt{g_W^2 + g_Y^2}
\end{equation}
Result: $m_Z = 91{,}188$~MeV ($<0.01\%$ from PDG). The Z mass
is the most precisely tested BCT electroweak prediction: agreement
to better than one part in ten thousand with no free parameters.

%----------------------------------------------------------------------
\section{Letter 16 --- The Higgs Boson Mass at Leading Order}

\medskip
The Higgs mass at leading order is derived from the BCT amplitude
mode of the condensate. The Higgs is identified as the radial
excitation of the superfluid order parameter. At leading order:
\begin{equation}
m_H^{(0)} = \sqrt{2\lambda}\,v
\end{equation}
where $\lambda = \alpha_0\pi$ is the BCT quartic coupling. The
leading result gives $m_H^{(0)} = 124{,}916$~MeV ($-0.27\%$).
The NLO correction is derived in Letter~134 (Volume~VII).

%----------------------------------------------------------------------
\section{Letter 17 --- The Fermi Constant}

\medskip
The Fermi constant is derived from the BCT weak coupling:
\begin{equation}
G_F = \frac{\pi\alpha}{\sqrt{2}\,m_W^2\sin^2\theta_W}
= \frac{1}{\sqrt{2}\,v^2}
\end{equation}
Result: $G_F = 1.16638 \times 10^{-5}$~GeV$^{-2}$ ($<0.01\%$).
The Fermi constant is determined entirely by the BCT electroweak
VEV, which is itself a geometric output. Zero free parameters.

%----------------------------------------------------------------------
\section{Letter 18 --- The BCT Uniqueness Theorem}
\noindent\textbf{DOI:} \href{https://doi.org/10.5281/zenodo.18957202}{10.5281/zenodo.18957202}

\medskip
The BCT Uniqueness Theorem proves that the BCT lattice at $c/a = \sqrt{2}$
is the \emph{only} crystalline vacuum state satisfying five simultaneous
physical requirements:
\begin{enumerate}
  \item Lorentz invariance (isotropic phonon propagation)
  \item Three fermion generations ($\mathbb{Z}_3$ outer automorphism)
  \item Maximum packing density in the parent dimension
  \item Modular invariance (self-dual lattice)
  \item The correct Standard Model gauge group
\end{enumerate}
No other lattice in any dimension satisfies all five. The vacuum has
no choice: if it crystallises, it must crystallise into the BCT lattice.

%----------------------------------------------------------------------
\section{Letter 19 --- The Electron Mass}
\noindent\textbf{DOI:} \href{https://doi.org/10.5281/zenodo.18905765}{10.5281/zenodo.18905765}

\medskip
The electron mass is derived from the BCT instanton action:
\begin{equation}
m_e = m_P\,\exp\!\left(-\frac{S_{D4}}{1-4\alpha_0/\pi}\right)
= 0.511000~\mathrm{MeV}
\end{equation}
Error: $<0.01\%$ from the PDG value. The electron mass is determined
by a single exponential suppression of the Planck mass by the D4
instanton action, corrected by the BCT self-consistency factor
$(1-4\alpha_0/\pi)$. No free parameters.

%----------------------------------------------------------------------
\section{Letter 20 --- General Relativity to All Orders}
\noindent\textbf{DOI:} \href{https://doi.org/10.5281/zenodo.18905870}{10.5281/zenodo.18905870}

\medskip
Full general relativity to all orders is derived from BCT condensate
geometry using the Jacobson--Parentani analogue gravity theorem.
The Einstein--Hilbert action emerges from the BCT stiff-condensate
condition $c_s = c$, which follows exactly from the Josephson coupling.
The Gauss--Bonnet coefficient $\alpha_{\mathrm{GB}} = \ell_P^2/(8\pi)$
is derived exactly, with no free parameters. GR is the unique
low-energy effective theory of the BCT vacuum.

%----------------------------------------------------------------------
\section{Letter 21 --- Muon and Tau Masses from C$_{4v}$ Geometry}
\noindent\textbf{DOI:} \href{https://doi.org/10.5281/zenodo.18906049}{10.5281/zenodo.18906049}

\medskip
The muon and tau masses are derived from BCT $C_{4v}$ one-loop geometry
via a Koide angle correction $\delta\theta_K = \varepsilon^2/(16\pi)$
where $\varepsilon = \sqrt{2}-1$. Results:
\begin{itemize}
  \item $m_\mu = 105.641$~MeV ($-0.016\%$)
  \item $m_\tau = 1776.894$~MeV ($+0.002\%$)
\end{itemize}
The correction reduces the muon error from $-1.60\%$ to $-0.016\%$:
a factor of 100 improvement from a single derived term with zero
free parameters.

%----------------------------------------------------------------------
\section{Letter 22 --- The QCD Confinement Scale}
\noindent\textbf{DOI:} \href{https://doi.org/10.5281/zenodo.18906377}{10.5281/zenodo.18906377}

\medskip
The QCD confinement scale is derived via the D4 theta-function
decomposition:
\begin{equation}
S_{\mathrm{QCD}} = \frac{\pi^5}{6}\!\left(1 - \rtet + \frac{\alpha_0}{2}\right)
= 45.4608
\end{equation}
giving $\LQCD = m_P\,e^{-S_{\mathrm{QCD}}} = 220.4$~MeV ($+0.20\%$,
Prediction~\#96). Zero free parameters.

%----------------------------------------------------------------------
\section{Letter 23 --- The Strong Coupling Constant}

\medskip
The strong coupling constant is derived from $\LQCD$:
\begin{equation}
\alpha_s(m_Z) = \frac{2\pi}{b_0\ln(m_Z/\LQCD)}
\end{equation}
with $b_0 = 23/3$ at one loop. BCT gives $\alpha_s(m_Z) = 0.1180$
($<0.01\%$ from PDG). The BCT derivation provides the first
first-principles derivation of $\alpha_s$ from pure geometry.

%----------------------------------------------------------------------
\section{Letter 24 --- The Electroweak Precision Observables}

\medskip
The complete set of electroweak precision observables is derived
and compared with LEP and LHC measurements. BCT achieves $<0.1\%$
agreement on 12 of 14 precision observables with zero free parameters.
The two residuals ($R_\tau$ at $-0.1\%$ and $\Delta(1232)$ at $-0.6\%$)
are identified as targets for higher-loop corrections in subsequent letters.

%----------------------------------------------------------------------
\section*{Volume II Summary}
%----------------------------------------------------------------------

\begin{table}[h]
\centering
\caption{Selected key results from Volume~II}
\begin{tabular}{lcc}
\toprule
Observable & BCT & Error \\
\midrule
$v_{\mathrm{EW}}$ (GeV) & 246.2198 & $-9.5\times10^{-5}\%$ \\
$m_W$ (MeV) & 80{,}370 & $+0.004\%$ \\
$m_Z$ (MeV) & 91{,}188 & $<0.01\%$ \\
$m_H$ (GeV) & 124.916 & $-0.27\%$ \\
$G_F$ ($10^{-5}$~GeV$^{-2}$) & 1.16638 & $<0.01\%$ \\
$m_e$ (MeV) & 0.511000 & $<0.01\%$ \\
$m_\mu$ (MeV) & 105.641 & $-0.016\%$ \\
$m_\tau$ (MeV) & 1776.894 & $+0.002\%$ \\
$\LQCD$ (MeV) & 220.4 & $+0.20\%$ \\
$\alpha_s(m_Z)$ & 0.1180 & $<0.01\%$ \\
\bottomrule
\end{tabular}
\end{table}

%----------------------------------------------------------------------
\begin{thebibliography}{9}

\bibitem{monograph}
M.~R.\ Cabri\'e,
\textit{The BCT Superfluid Lattice Model: Complete Derivation of
Standard Model Parameters from Vacuum Geometry},
Zenodo (2026),
\href{https://doi.org/10.5281/zenodo.18884976}{10.5281/zenodo.18884976}.

\bibitem{appvol1}
M.~R.\ Cabri\'e,
\textit{Supplementary Material Volume 1},
Zenodo (2026),
\href{https://doi.org/10.5281/zenodo.18885133}{10.5281/zenodo.18885133}.

\bibitem{appvol2}
M.~R.\ Cabri\'e,
\textit{Supplementary Material Volume 2},
Zenodo (2026),
\href{https://doi.org/10.5281/zenodo.18885472}{10.5281/zenodo.18885472}.

\bibitem{strongCP}
M.~R.\ Cabri\'e,
\textit{Geometric Resolution of the Strong CP Problem without an Axion},
Zenodo (2026),
\href{https://doi.org/10.5281/zenodo.18885628}{10.5281/zenodo.18885628}.

\bibitem{vol1}
M.~R.\ Cabri\'e,
\textit{BCT Letters Collection --- Volume I: Foundations},
Zenodo (2026).

\end{thebibliography}

\end{document}
